% ---
% id: reductions/cor:monotone
% kind: cor
% title: "Monotonicity of the worst-case criterion along the pre-order"
% summary: >-
%   If one target squashes to a cleaner one, its worst-case criterion over a
%   family of inputs equals the cleaner target's over the family's image under
%   the connecting channel, and is at most the cleaner target's over any larger
%   family. Moving up the pre-order can only raise the worst case.
% prominence: standard
% status: accepted
% uses: [thm:factorization]
% assumes: []
% refs: []
% provenance: newly-derived
% lean:
% ---

Composing squashers costs nothing in soundness; the corollary says in which
direction it costs tightness. Moving \emph{up} the pre-order, to a cleaner
target, hands the adversary more.

\begin{corollary}\label{cor:monotone}
For a POVM $G$ and a family $S$ of inputs write
$\varepsilon^{*}(G,S)\coloneqq\sup_E\sup_{\sigma\in S_E}
\Crit(\Eprot_G,\sigma)$. Let $F_2\dsq_{\Lambda_{12}}F_1$ and let both
\tref[squashed protocol]{squashed protocols} use the same processing $\bar{\Grnd}_j$ and $\Kfin$. Then for
every family $S$,
\begin{equation*}
  \varepsilon^{*}(F_2,S)
  =\varepsilon^{*}\big(F_1,(\idmap\otimes\Lambda_{12}^{\otimes n}\otimes\idmap_E)(S)\big)
  \;\le\;\varepsilon^{*}(F_1,S_1)
  \quad\text{for every }S_1\supseteq
  (\idmap\otimes\Lambda_{12}^{\otimes n}\otimes\idmap_E)(S).
\end{equation*}
\end{corollary}

\begin{proof}
The equality is Theorem~\ref{thm:factorization} applied with $(F_2,F_1)$ in
place of $(M,F)$, pointwise in $\sigma$ and then in the suprema; the inequality
is monotonicity of a supremum in the set.
\end{proof}
